% Standalone problem document, generated from the adjacent README.md.
% Compile with XeLaTeX. Mathematical content is maintained in Markdown.
\documentclass[11pt,a4paper]{article}
\usepackage[fontsize=10.5pt]{scrextend}
\usepackage[margin=22mm,headheight=15pt,headsep=9mm,footskip=11mm]{geometry}
\usepackage{fontspec}
\setmainfont{texgyrepagella}[Extension=.otf,UprightFont=*-regular,BoldFont=*-bold,ItalicFont=*-italic,BoldItalicFont=*-bolditalic]
\setsansfont{texgyreheros}[Extension=.otf,UprightFont=*-regular,BoldFont=*-bold,ItalicFont=*-italic,BoldItalicFont=*-bolditalic]
\setmonofont{texgyrecursor}[Extension=.otf,UprightFont=*-regular,BoldFont=*-bold,ItalicFont=*-italic,BoldItalicFont=*-bolditalic,Scale=MatchLowercase]
\usepackage{amsmath,amssymb}
\usepackage{unicode-math}
\setmathfont{texgyrepagella-math.otf}
\usepackage{xcolor,microtype,enumitem,needspace,fancyhdr}
\usepackage{bookmark}
\definecolor{ink}{HTML}{17344B}
\definecolor{accent}{HTML}{197D80}
\definecolor{muted}{HTML}{596773}
\hypersetup{colorlinks=true,linkcolor=accent,urlcolor=accent,pdftitle={AC-06 - Explicit
tensors with quadratic border
rank},pdfauthor={Open Problems in Numerical Linear Algebra}}
\urlstyle{same}
\setlength{\parindent}{0pt}
\setlength{\parskip}{5pt plus 1pt minus 1pt}
\linespread{1.04}
\setlength{\emergencystretch}{3em}
\widowpenalty=10000
\clubpenalty=10000
\displaywidowpenalty=10000
\allowdisplaybreaks[1]
\setlist{nosep,leftmargin=1.5em}
\providecommand{\tightlist}{\setlength{\itemsep}{0pt}\setlength{\parskip}{0pt}}
\providecommand{\pandocbounded}[1]{#1}
\pagestyle{fancy}
\fancyhf{}
\fancyhead[L]{\sffamily\footnotesize\color{muted}OPEN PROBLEMS IN NUMERICAL LINEAR ALGEBRA}
\fancyhead[R]{\sffamily\bfseries\footnotesize\color{accent}AC-06}
\fancyfoot[L]{\parbox[b]{0.9\textwidth}{\sffamily\fontsize{8}{10}\selectfont\color{muted}Editorial ratings; a bounded search cannot certify that a problem remains unsolved.\\Literature check: 2026-09-08}}
\fancyfoot[R]{\sffamily\footnotesize\color{muted}\thepage}
\renewcommand{\headrulewidth}{0.3pt}
\renewcommand{\headrule}{\hbox to\headwidth{\color{accent}\leaders\hrule height \headrulewidth\hfill}}
\makeatletter
\renewcommand{\section}{\Needspace{5\baselineskip}\@startsection{section}{1}{0pt}{12pt plus 2pt minus 2pt}{4pt}{\sffamily\bfseries\large\color{ink}}}
\renewcommand{\subsection}{\Needspace{4\baselineskip}\@startsection{subsection}{2}{0pt}{9pt plus 1pt minus 1pt}{3pt}{\sffamily\bfseries\normalsize\color{ink}}}
\makeatother
\setcounter{secnumdepth}{0}
\begin{document}
{\sffamily\small\color{muted}Arithmetic and complexity\par}
\vspace{3pt}
{\raggedright\hyphenpenalty=10000\exhyphenpenalty=10000\sffamily\bfseries\fontsize{21}{24}\selectfont\color{ink}Explicit
tensors with quadratic border rank\par}
\vspace{7pt}
{\sffamily\small\textbf{Difficulty:} extreme\quad\textbf{Importance:} broadly
interesting\par}
\vspace{4pt}
{\color{accent}\hrule height 0.8pt}
\vspace{6pt}
\textbf{Topic:} explicit lower bounds for tensor decompositions

\textbf{Status:} open; admitted after a bounded literature search found
no resolution.

\subsection{Context and notation}\label{context-and-notation}

Over \(\mathbb C\), the tensor rank \(R(T)\) is the minimum number of
pure tensors in an exact sum for \(T\). The border rank
\(\underline R_{\mathbb C}(T)\) is the least \(r\) such that \(T\) is a
Euclidean limit of complex tensors of rank at most \(r\).

\subsection{Problem statement}\label{problem-statement}

Construct a deterministic algorithm that, on input the positive integer
\(n\), outputs the rational entries of a tensor
\(T_n\in\mathbb Q^{n\times n\times n}\) in time polynomial in \(n\), and
prove that there are constants \(c>0,n_0\) such that
\(\underline R_{\mathbb C}(T_n)\ge c n^2\) for all \(n\ge n_0\). Output
rationals are encoded by binary integer numerators and denominators;
thus the time bound also limits their bit lengths. Border rank is over
\(\mathbb C\), as defined in the notation.

\subsection{Why it matters}\label{why-it-matters}

Explicit hard tensors would supply lower bounds for bilinear computation
that are much stronger than the currently available examples. Generic
existence does not provide the required efficient construction.

\subsection{References and status}\label{references-and-status}

M. Michałek and J. Landsberg,
\href{https://www.theoryofcomputing.org/articles/v021a013/v021a013.pdf}{\emph{Towards
Finding Hay in a Haystack: Explicit Tensors of Border Rank Greater Than
2.02m}}, Theory of Computing 21(13), 2025, §1 and Definition 1.1,
distinguishes polynomial-time explicitness from weaker notions and
obtains a linear bound. The quadratic target and rational-output
encoding are an editorial quantitative formulation of the paper's
explicit-versus-generic challenge, not a numbered conjecture quoted from
it. Searches for ``explicit tensors superlinear rank 2026'' and
``explicit tensors superlinear border rank'' located no quadratic
construction; weaker semi-explicit constructions do not meet this
statement. \textbf{Admitted: no resolution located.}
\end{document}
